\documentclass[10pt,a4paper]{amsart}
\usepackage[margin=19mm]{geometry}
\usepackage{amsmath,amssymb,mathtools}
\usepackage[scr=rsfs,cal=euler]{mathalfa}
\usepackage{xfrac}
\usepackage[unicode,hidelinks,bookmarks=false]{hyperref}
\newcommand{\ie}{i.e.\ }
\newcommand{\cf}{cf.\ }
\newcommand{\resp}{respectively\ }
\newcommand{\Subsection}{\subsection}
\providecommand{\nobreakdash}{\nobreak\hbox{-}}
\setlength{\emergencystretch}{2em}
\multlinegap0pt
\usepackage{amssymb}
\usepackage{mathtools}
\usepackage{xcolor}
\usepackage{esint}
\usepackage[shortlabels]{enumitem}
\newtheorem{thm}{Theorem}
\newtheorem{proposition}{Proposition}
\newcommand{\R}{\mathsf{R}}
\newcommand{\RCD}{\mathsf{RCD}}
\newcommand{\Ric}{\mathsf{Ric}}
\newcommand{\bb}[1]{\mathbb{#1}}
\newcommand{\de}{{\rm d}}
\newcommand{\m}{\mathfrak{m}}
\newcommand{\sd}{\mathsf{d}}
\newcommand{\ssf}[1]{\mathsf{#1}}
\title{On manifolds with almost non-negative Ricci curvature and integrally-positive $k^{th}$-scalar curvature}
\author{Alessandro Cucinotta, Andrea Mondino}
\date{Source: arXiv:2504.06865v3 (CC BY 4.0)}
\begin{document}
\maketitle

\noindent\emph{Source and licence.} On manifolds with almost non-negative Ricci curvature and integrally-positive $k^{th}$-scalar curvature. Version arXiv:2504.06865v3 (CC BY 4.0), distributed by arXiv under the Creative Commons Attribution 4.0 International licence, http://creativecommons.org/licenses/by/4.0/, as recorded in the arXiv licence metadata for this exact version and shown on the abstract page https://arxiv.org/abs/2504.06865v3. Full licence notice: \texttt{licenses/arxiv-2504.06865-CC-BY-4.0.txt}.\par\smallskip

\noindent\textit{Editorial scope.} Counted unit: the article's thm CT2 (source0.tex lines 939--955). The article's complete original proof of it is delivered with it (source0.tex lines 957--1033, one \texttt{proof} environment of 77 lines). No mathematics is written, repaired or re-derived: the statement and the proof are the article's own lines, in the article's own notation, and no retained line is substituted or re-flowed.  Retained uncounted as the source-local context the proof needs: lines 180--190; lines 473--475; lines 490--492.  Nothing was omitted from any retained span, so the package carries no omission record.  No retained line contains a citation, so the wrapper carries no bibliography and no bibliography entry.  No retained line contains a figure environment, an image-inclusion command or drawing code, so no image is delivered and the package declares no asset dependency.  Editorial formatting only, in addition: an AMS \texttt{amsart} preamble that restates the article's own declarations and macros used by the retained lines, namely the environments \texttt{thm}, \texttt{proposition} on separate counters, together with the standard packages those lines need.  The retained environments print in this document as Theorem 1, Proposition 1 in order of appearance, whereas the article prints the same items with the numbering of its own full document.  The complete native source of the article as distributed in the arXiv e-print for this exact version is bundled unchanged as \texttt{sources/176246/source0.tex}.
\begin{thm} \label{CT1}
    Let $L,k,n \in \bb{N}$ with $k \leq n$ and $s \in (0,1)$ be fixed. Then 
    \begin{align*}
    &\sup \Big\{\fint_{B_1(x)} \R_{k} \wedge L \, \de \ssf{Vol} : (M,g,x) \text{ has } \Ric_{M} \geq -\delta, \text{ and } B_{10}(x) \text{ is  $(k,\delta)$-symmetric} 
    \Big\} ,
    \\
    &  \sup \Big\{\fint_{B_1(x)} |\R|^s \, \de \ssf{Vol} : (M,g,x) \text{ has } \Ric_{M} \geq -\delta, \text{ and } B_{10}(x) \text{ is } \delta \text{-regular} 
    \Big\}  
    \end{align*}
    converge to $0$ as $\delta \searrow 0$.
    \end{thm}

\begin{thm} \label{T2}
    Let $(X, \sd, \bar{x})$ be a non-collapsed Ricci limit space of dimension $n \geq 2$. Then $\mathcal{S}_{n-1} \setminus \mathcal{S}_{n-2} = \emptyset$.
\end{thm}

\begin{proposition} \label{P1}
    Let $(X,\sd,\m)$ be an $\RCD(K,N)$ space of essential dimension equal to $1$. Then $(X, \sd)$ is isometric to a closed connected subset of $\bb{R}$ or to $\mathbb{S}^1(r):=\{x \in \bb{R}^2: |x|=r \}$.
\end{proposition}

\begin{thm} \label{CT2}
    Let $s \in (0,1)$, let $\epsilon,v,L \in (0,+\infty)$ and let $k,n \in \bb{N}$ with $k \leq n$. 
    \begin{enumerate}
        \item \label{Thm4.7-1} There exists $\delta>0$ such that the following holds. Let $(M^n,g,p)$ be a manifold with $\Ric_{M} \geq -\delta$ on $M$ and
    \[
     \fint_{B_1(p)} |\R_{M}|^s \, \de \ssf{Vol}  \geq \epsilon, \quad \ssf{Vol}(B_1(p)) \geq v.
    \] 
    Then $\sd_{\rm{pGH}}(M,Y) \geq \delta$, for any metric space $(Y,\sd_y,y)$  splitting $\bb{R}^{n-1}$ isometrically.
    \item \label{Thm4.7-2}
    There exists $\delta>0$ such that the following holds. Let $(M^n,g,p)$ be a manifold with $\Ric_{M} \geq -\delta$ on $M$, and
    \[
    \fint_{B_1(p)} \ssf{Ric}_{M,n-2} \wedge 0 \, \de \ssf{Vol}  \geq -\delta, \quad 
     \fint_{B_1(p)} \R_{M} \wedge L \, \de \ssf{Vol}  \geq \epsilon.
    \] 
     Then $\sd_{\rm{pGH}}(M,Y) \geq \delta$, for any metric space $(Y,\sd_y,y)$  splitting $\bb{R}^{n-1}$ isometrically.
    \end{enumerate}
\end{thm}

\begin{proof} 
\textbf{Proof of \ref{Thm4.7-1}.}
Assume by contradiction that there exist $\epsilon_0>0$, $v_0>0$, $\delta_j \downarrow 0$ and pointed Riemannian manifolds $(M_j,g_j,p_j)$ with $\Ric_{M_j} \geq -\delta_j$, $\ssf{Vol}_j(B_1(p_j)) \geq v_0$, and
\[
\fint_{B_1(p_j)}|\R_{M_j}|^s \, \de \ssf{Vol}_j \geq \epsilon_0,
\]
such that
$M_j \to X$ in  pGH-sense, where $(X,\sd,p)$ splits $\bb{R}^{n-1}$ isometrically.
Thanks to Proposition \ref{P1} and Theorem \ref{T2}, then either $X=\bb{R}^n$ or $X=\bb{R}^{n-1} \times S^1_r$, for some $r>0$. 
\smallskip

\textbf{Case $X=\bb{R}^n$}.
Applying the second assertion of Theorem \ref{CT1} to $(M_j, g_j, p_j)$, for $j$ large enough, we reach a contradiction.

\smallskip
\textbf{Case $X=\bb{R}^{n-1} \times S^1_r$}.  
Let $s \in (0,1)$. We claim that for every $j$ there exists $p_j^s \in B_1(p_j)$ such that
\begin{equation} \label{eqcorretta}
\fint_{B_s(p_j^s)} |\R_{M_j}|^s \, d\ssf{Vol}_j \geq c(n,s)\epsilon_0.
\end{equation}
If \eqref{eqcorretta} holds, we reach a contradiction by applying the second assertion of Theorem \ref{CT1} to the rescaled sequence $(M_j, s^{-2} g_j, p_j)$, for $s \in (0,1)$ small enough independent of $j$.

We prove \eqref{eqcorretta}. Let $j$ be fixed. Let $\{q_i\}_{i=1}^l \subset B_1(p_j)$ be such that $\{B_s(q_i)\}_{i=1}^l$ covers $B_1(p_j)$, while the balls $\{B_{s/5}(q_i)\}_{i=1}^l$ are disjoint. Assume without loss of generality that
\[
\int_{B_s(q_1)}|\R_{M_j}|^s \, d\ssf{Vol}_j=\max_{i}\int_{B_s(q_i)}|\R_{M_j}|^s \, d\ssf{Vol}_j.
\]
Then,
\[
\int_{B_1(p_j)}|\R_{M_j}|^s \, d\ssf{Vol}_j \leq l \int_{B_s(q_1)}|\R_{M_j}|^s \, d\ssf{Vol}_j.
\]
Hence, to conclude, it suffices to show that $l \leq c(n,s)\frac{\ssf{Vol}_j(B_1(p_j))}{\ssf{Vol}_j(B_s(q_1))}$.
To this aim, note that
\begin{align*}
\ssf{Vol}_j(B_1(p_j))  \geq c(n)\ssf{Vol}_j(B_2(p_j)) \geq c(n)\sum_{i=1}^l \ssf{Vol}_j(B_{s/5}(q_i)) \\
 \geq c(n,s)l\ssf{Vol}_j(B_s(q_1)).
\end{align*}
This concludes the proof of item \ref{Thm4.7-1}.
\smallskip

\textbf{Proof of \ref{Thm4.7-2}.}
By contradiction, assume there exist numbers $\delta_j \downarrow 0$ and pointed Riemannian manifolds $(M_j,g_j,p_j)$ with $\Ric_{M_j} \geq -\delta_j$ on $M_j$, and
\begin{equation} \label{E4}
\fint_{B_1(p_j)} \ssf{Ric}_{M_j,n-2} \wedge 0 \, \de \ssf{Vol}_j  \geq -\delta_j, \quad 
\fint_{B_1(p_j)}\R_{M_j} \wedge L \, \de \ssf{Vol}_j \geq \epsilon,
\end{equation}
such that
$M_j \to X$ in pGH-sense, where $(X,\sd,p)$ splits $\bb{R}^{n-1}$ isometrically. 

By the first assertion of Theorem \ref{CT1}, it holds
\begin{equation}\label{eq:Rjn-1to0}
\fint_{B_1(p_j)} \R_{M_j,n-1} \wedge L \,\de\ssf{Vol}_j \to 0, \quad \text{ as } j\to \infty.
\end{equation}
Consider an orthonormal base $\{e_1,\cdots,e_n\}$ for the tangent space of the manifold $M_j$ in a point $m \in M_j$, and assume that $\Ric_{M_j}$ in this basis is represented by a diagonal matrix $A=(a_{h,l}) \in \bb{R}^{n \times n}$ such that $a_{h,h} \leq a_{h+1,h+1}$ for every $h \in \{1, \cdots ,n-1\}$. It holds
\begin{align*}
\R_{M_j} = \sum_{h \neq l} \ssf{Sec}_{M_j}(e_h,e_l)&=\R_{M_j,n-1}+\sum_{h=1}^{n-1} \ssf{Sec}_{M_j}(e_n,e_h) \\
&=2\R_{M_j,n-1}-\sum_{\substack{ h \neq n \\ l \neq n}}\ssf{Sec}_{M_j}(e_h,e_l).
\end{align*}
Hence,
\begin{equation} \label{eqff}
\R_{M_j} \wedge L \leq 2(0 \vee \R_{M_j,n-1} \wedge L)-\sum_{h=1}^{n-1}\Big(\sum_{\substack{ l \neq h \\ l \neq n}}\ssf{Sec}_{M_j} (e_h,e_l)\Big) \wedge 0. 
% \\
% & \leq  2(\R_{M_j,n-1} \wedge L)+2n\delta_j-n\ssf{Ric}_{M_j,n-2} \wedge 0.
\end{equation}
By the first inequality of \eqref{E4}, for every $h \in \{1,\cdots,n-1\}$, it follows
\[
\fint_{B_1(p_j)}\Big(\sum_{\substack{ l \neq h \\ l \neq n}}\ssf{Sec}_{M_j} (e_h,e_l)\Big) \wedge 0 \, d\ssf{Vol} \geq -\delta_j.
\]
Combining with \eqref{eqff}, it holds
\[
\fint_{B_1(p_j)} \R_{M_j} \wedge L \,\de\ssf{Vol}_j \leq 2\fint_{B_1(p_j)} \R_{M_j,n-1} \wedge L \,\de\ssf{Vol}_j + c(n)\delta_j.
\]
From \eqref{eq:Rjn-1to0}, we infer that
\[
\limsup_{j \to +\infty} \fint_{B_1(p_j)} \R_{M_j} \wedge L \,\de\ssf{Vol}_j= 0,
\]
contradicting the second inequality of \eqref{E4}.
\end{proof}

\end{document}
